\section{Experiment: Exploring Compilation Internals}

\subsection{Objective}
Observe the compiler's transformation from source code to assembly, including intermediate representations and optimization effects.

\subsection{Setup Requirements}
\begin{itemize}
\item GCC compiler
\item LLVM/Clang (optional, for LLVM IR)
\item Linux environment
\end{itemize}

\subsection{Part 1: Generating Assembly Output}

Create a simple C file:

\begin{lstlisting}[language=C]
// simple.c
int add(int a, int b) {
    return a + b;
}

int main(void) {
    int result = add(3, 5);
    return result;
}
\end{lstlisting}

Generate assembly without optimization:

\begin{lstlisting}[language=bash]
gcc -S -O0 simple.c -o simple_O0.s
\end{lstlisting}

Generate assembly with optimization:

\begin{lstlisting}[language=bash]
gcc -S -O2 simple.c -o simple_O2.s
gcc -S -O3 simple.c -o simple_O3.s
\end{lstlisting}

Compare the outputs:

\begin{lstlisting}[language=bash]
diff simple_O0.s simple_O2.s
\end{lstlisting}

\subsection{Part 2: Viewing GCC Intermediate Representations}

Create a more complex function:

\begin{lstlisting}[language=C]
// compute.c
int compute(int n) {
    int sum = 0;
    for (int i = 1; i <= n; i++) {
        sum += i * i;
    }
    return sum;
}
\end{lstlisting}

Dump GIMPLE representation:

\begin{lstlisting}[language=bash]
gcc -fdump-tree-gimple compute.c -c -o compute.o
cat compute.c.003t.gimple
\end{lstlisting}

View all optimization passes:

\begin{lstlisting}[language=bash]
gcc -fdump-tree-all -O2 compute.c -c -o compute.o
ls -la compute.c.*.tree
\end{lstlisting}

\subsection{Part 3: Viewing LLVM IR}

Generate LLVM IR:

\begin{lstlisting}[language=bash]
clang -S -emit-llvm compute.c -o compute.ll
cat compute.ll
\end{lstlisting}

Generate optimized IR:

\begin{lstlisting}[language=bash]
clang -S -emit-llvm -O2 compute.c -o compute_O2.ll
diff compute.ll compute_O2.ll
\end{lstlisting}

\subsection{Part 4: Optimization Effects}

Create test code:

\begin{lstlisting}[language=C]
// opt_test.c
int dead_code(int x) {
    int a = 10;
    int b = 20;
    int c = a + b;  // Dead: never used
    return x * 2;
}

int constant_fold(void) {
    int x = 2 + 3 * 4;
    return x;
}
\end{lstlisting}

Generate assembly at different optimization levels:

\begin{lstlisting}[language=bash]
gcc -S -O0 opt_test.c -o opt_O0.s
gcc -S -O2 opt_test.c -o opt_O2.s
\end{lstlisting}

Compare dead code elimination:

\begin{lstlisting}[language=bash]
grep -A20 "dead_code:" opt_O0.s
grep -A10 "dead_code:" opt_O2.s
\end{lstlisting}

\subsection{Part 5: Inlining Analysis}

Create test file:

\begin{lstlisting}[language=C]
// inline_test.c
static int square(int x) {
    return x * x;
}

int calculate(int n) {
    return square(n) + square(n + 1);
}
\end{lstlisting}

Examine with and without inlining:

\begin{lstlisting}[language=bash]
gcc -S -O0 -fno-inline inline_test.c -o inline_no.s
gcc -S -O2 inline_test.c -o inline_yes.s
\end{lstlisting}

\subsection{Questions for Reader}

\begin{enumerate}
\item What differences do you observe between \texttt{-O0} and \texttt{-O2} for the \texttt{add} function?
\item In the SSA dump, look for phi functions. Where do they appear and what do they represent?
\item How much code was eliminated from \texttt{dead\_code} at \texttt{-O2}?
\item Compare the loop in \texttt{loop\_optimize} at different optimization levels. What transformations do you see?
\item What are the advantages and disadvantages of function inlining?
\end{enumerate}

\subsection{Further Exploration}
\begin{itemize}
\item Explore GCC's \texttt{-fverbose-asm} flag for annotated assembly
\item Use \texttt{-fopt-info} to see what optimizations the compiler applies
\item Compare GCC and Clang optimization strategies
\item Study the \texttt{-march=native} flag for CPU-specific optimizations
\end{itemize}
